% Working draft explicitly requested by Adam Topaz, 9 November 2026.
% Standalone: latexmk -pdf cohomological_hyperplane_gt_blueprint.tex
\documentclass[11pt]{article}
\usepackage[T1]{fontenc}
\usepackage{lmodern,amsmath,amssymb,amsthm,mathtools,mathrsfs}
\usepackage[margin=1in]{geometry}
\usepackage{microtype,enumitem,booktabs,longtable,array,xcolor,xurl,listings,needspace}
\usepackage[colorlinks=true,linkcolor=blue!45!black,urlcolor=blue!45!black]{hyperref}
\usepackage{fancyhdr}
\pagestyle{fancy}\fancyhf{}
\fancyhead[L]{\small Cohomological hyperplane GT: formalization blueprint}
\fancyhead[R]{\small\thepage}
\setlength{\headheight}{14pt}
\setlength{\emergencystretch}{3em}
\setlist{itemsep=3pt,topsep=5pt}
\lstset{basicstyle=\ttfamily\small,breaklines=true,columns=fullflexible,keepspaces=true,frame=single}
\newtheorem{theorem}{Theorem}[section]
\newtheorem{lemma}[theorem]{Lemma}
\newtheorem{proposition}[theorem]{Proposition}
\newtheorem{definition}[theorem]{Definition}
\newtheorem{corollary}[theorem]{Corollary}
\theoremstyle{remark}\newtheorem{remark}[theorem]{Remark}
\newcommand{\A}{\mathbb A}
\newcommand{\Pj}{\mathbb P}
\newcommand{\Gm}{\mathbb G_m}
\newcommand{\Z}{\mathbb Z}
\newcommand{\Q}{\mathbb Q}
\newcommand{\C}{\mathbb C}
\newcommand{\F}{\mathbb F}
\newcommand{\cH}{\mathcal H}
\newcommand{\cD}{\mathcal D}
\newcommand{\OS}{\ensuremath{\operatorname{OS}}}
\newcommand{\Aut}{\operatorname{Aut}}
\newcommand{\Gal}{\operatorname{Gal}}
\newcommand{\Hom}{\operatorname{Hom}}
\newcommand{\Ann}{\operatorname{Ann}}
\newcommand{\Span}{\operatorname{span}}
\newcommand{\im}{\operatorname{im}}
\newcommand{\codim}{\operatorname{codim}}
\newcommand{\id}{\operatorname{id}}
\newcommand{\dR}{\mathrm{dR}}
\newcommand{\code}[1]{\nolinkurl{#1}}
\newcommand{\node}[3]{\par\Needspace{8\baselineskip}\smallskip\noindent\textbf{Node #1.}\quad
 \textsf{Proposed:}\ \code{#2}\par\noindent
 \textsf{Dependencies:} #3.\par\smallskip}
\title{Cohomological hyperplane reconstruction\\
\large A detailed Lean blueprint from the existing Orlik--Solomon library}
\author{Working draft prepared for Adam Topaz}
\date{9 November 2026}
\begin{document}
\maketitle
\begin{abstract}
We specify a formal proof that the automorphism group of the marked
degree-one and degree-two cohomology diagram of hyperplane complements is
$R^\times\times\Gal(k/k_0)$, for a nonzero local coefficient ring or an
integral domain. Every new mathematical step is proved below. The only
imported mathematical results are existing theorems of this repository
or the pinned Mathlib. In particular the local reconstruction is proved
by elementary equations, and projective reconstruction is proved in
coordinates rather than delegated to an unformalized fundamental theorem.
The output includes the C1--C5 theorem and an unconditional specialization
to the constructed algebraic de Rham theory. Proposed Lean names and
signatures are specifications, not elaborated declarations.
\end{abstract}
\tableofcontents
\newpage

\Needspace{14\baselineskip}
\section{Exact target, scope, and existing inputs}
\label{sec:boundary}

This document formalizes the new proof in
\code{linear_gt_draft/cohomological_hyperplane_gt.tex}. It does not require
that manuscript or the Pop--Topaz paper as an unproved mathematical input:
the new proofs are repeated here, with their interfaces and intermediate
claims. Source locations refer to repository commit \code{1d89e7d}
and Mathlib \code{6db5db0d2279846cc8fb25901bd53221d37a77d8}.
The existing architecture remains authoritative for existing names.
All modules under \code{OrlikSolomon/HyperplaneGT/} mentioned here are
proposed additions; this draft changes neither Lean code nor architecture.

Let $k_0\subset k$ be a perfect field and its algebraic closure, let
$n\geq2$, and let $S\subset k_0$ contain $0$ and generate $k_0$ over
the prime field. Let $\mathscr F$ be a cofinal family of finite
$\Gal(k/k_0)$-stable arrangements in $\A^n_k$. Put
$T=\A^1_k\setminus\{0,1\}$. The diagram has complements $U_A$,
inclusions, and maps to the same object $T$ given by $x_i-s$ and
$x_i-x_j$ ($s\in S$, $i<j$), whenever their two excluded fibers are
in $A$. The reconstruction theorem is
\[
 R^\times\times(k\simeq_{k_0}k)
 \simeq\Aut_2(E|\cD)\simeq\Aut_{\rm gr}(E|\cD),
 \qquad (u,\sigma)e_H=u e_{\sigma H}.
\]
The first group consists of natural linear equivalences of degree one
preserving the full kernel of its exterior product into degree two.
The second consists of natural graded unital $R$-algebra equivalences
of the diagonal $\bigoplus_{q\geq0}E^{q,q}$. Scalars act in degree $q$
by $u^q$. Quotienting either group by this central subgroup gives
$\Gal(k/k_0)$. These are abstract group equivalences.

The purely algebraic reconstruction is stated first for the explicit
Orlik--Solomon diagram. It assumes no cohomology theory. An existing
structure theorem then identifies a supplied theory's diagram with
this explicit diagram. A new structure with a field asserting the
Orlik--Solomon theorem is neither needed nor permitted.

\subsection{Cohomology hypotheses, completely specified}

The supplied value of
\code{OrlikSolomon.Geometric.OrdinaryTheory}, with parameters
$D$ and $\mathcal B$, has the following
data, where $D$ may be the smooth affine domain or a larger admissible
domain, $R$ is commutative, and $B=\bigoplus_{w\in\Z}B_w$ is a
commutative $R$-algebra with $B_0\simeq R$ through the structural map.
\begin{enumerate}[label=\textbf{C\arabic*.},leftmargin=*]
\item Each $C(X)$ is a unital, possibly noncommutative $B$-algebra
internally bigraded by $R$-submodules $C(X)_{p,w}$, with central
coefficients $B_w$ in bidegree $(0,w)$ and coherent scalar towers.
\item Pullbacks along all morphisms in $D$ are functorial unital
$B$-algebra maps preserving degrees. Products of homogeneous classes
supercommute with the sign determined by the cohomological degrees.
\item $B\to C(\operatorname{Spec}k)$ is bijective, with the point
supported in cohomological degree zero, and
$C(X)\to C(X\times\A^1)$ is bijective.
\item For each certified principal smooth divisor pair
$D'\xrightarrow{i}X\xleftarrow{j}U$ there are $B$-linear Gysin and
residue maps of shifts $(2,1)$ and $(-1,-1)$ with
$\im g=\ker j^*$, $\im j^*=\ker r$, $\im r=\ker g$.
Here a certificate means a globally regular non-zero-divisor $f$,
the closed scheme $Z(f)$, its open complement $X_f$, and smoothness
of the three schemes.
\item Both maps commute with the specified morphisms of pairs, for
which the ambient pullback of the target equation is a unit times
the source equation. The right projection formula is
$r(\alpha j^*\beta)=r(\alpha)i^*\beta$ for
$\alpha\in C(U)_{1,1}$ and arbitrary $\beta\in C(X)$.
\end{enumerate}
These are existing inputs, with all their consequences already proved
for finite arrangements. They do not include logarithmic generation,
freeness, a dimension bound, residue surjectivity, Gysin vanishing, or
a short exact deletion sequence. They do not include a semilinear
action of a field automorphism. The diagonal must not be confused
with all weight-zero groups.

The theorem is conditional on this package in exactly the same sense
as the existing arrangement theorem. The de Rham package in
characteristic zero is already constructed and is an unconditional
specialization. Singular and geometric \'etale cohomology are examples
in the paper; their foundations and C1--C5 implementations are not
present in this checkout. Accordingly no named singular or \'etale
Lean instance is silently assumed in the dependency graph below.
Their applications of the new theorem are conditional on a value of
this existing interface. This is a complete blueprint for the new
axiomatic reconstruction proof, not a claim to have formalized the
foundations of every named cohomology theory.

\subsection{Verified repository entry points}

All names in this table are existing names, not proposals. Paths in
the first column are under \code{OrlikSolomon/}.
\small
\begin{longtable}{@{}p{.40\textwidth}p{.55\textwidth}@{}}
\toprule Module and location & Existing declarations used\\\midrule\endhead
\code{Cohomology/Geometric/Axioms.lean:49}&
\code{Geometric.OrdinaryTheory}; its specified scalar wrappers
\code{mapOverR}, \code{gysinOverR}, \code{residueOverR}.\\
\code{Cohomology/Axioms.lean:156}&
\code{OrdinaryTheory}, the equivalent algebraic interface.\\
\code{Presentation/Geometric/Ordinary.lean:192}&
\code{Geometric.OrdinaryTheory.degreeEquiv}; at line 221,
\code{diagonalEquiv}; at line 253,
\code{diagonalEquiv_generator_of_isEquation}.\\
\code{Presentation/Geometric/Affine.lean:332}&
\code{Geometric.OrdinaryTheory.affineEquiv_map}; at line 353,
\code{affineEquiv_comapEquiv}.\\
\code{Presentation/Basic.lean:71}&
\code{osRelations}; line 89, \code{OSAlgebra}; line 129,
\code{OSAlgebra.quotient_eq_zero_iff}.\\
\code{Presentation/Basic.lean:176}&
\code{OSAlgebra.generator_mul_eq_zero_of_disjoint}; line 183,
\code{wordBoundary_eq_zero_of_isDependent}.\\
\code{Presentation/Basic.lean:348}&
\code{OSAlgebra.map}, \code{map_generator}, \code{map_comp};
line 397, \code{degree}; line 430, \code{span_word_length}.\\
\code{Presentation/Generators.lean:149}&
\code{OSAlgebra.linearIndependent_generator}.\\
\code{Algebra/RelationQuotient.lean:101}&
\code{OrlikSolomon.relationSubmodule}; line 144,
\code{relationSubmodule_le_of_span_eq_top}; line 238,
\code{relationSubmodule_isHomogeneous}.\\
\code{Geometry/Hyperplane.lean:227}&
\code{IsHyperplane.exists_isEquation}; line 287,
\code{IsEquation.existsUnique_smul}; line 455,
\code{IsHyperplane.codim_inf_eq_two}.\\
\code{Geometry/Hyperplane.lean:513}&
\code{codim_iInf_of_isEquation}, the rank formula for nonempty
intersections.\\
\code{Geometry/Complement.lean:468}&
\code{SmoothAlg.complementMap}, arbitrary permitted affine pullback
of complement rings.\\
\code{Cohomology/Logarithms.lean:315}&
\code{OrdinaryTheory.map_logClass}; the geometric version is
\code{Geometric.OrdinaryTheory.pull_logClass} in
\code{Cohomology/Geometric/Logarithms.lean:309}.\\
\code{Presentation/BaseChange.lean:116}&
\code{OSAlgebra.baseChangeEquiv}, with its generator formula.\\
\code{DeRham/Schemes/Instance.lean:104}&
\code{Geometric.DeRham.ordinaryTheory}.\\
\code{DeRham/Schemes/Corollaries.lean:81}&
\code{Geometric.DeRham.affineEquiv}; line 124,
\code{affineEquiv_generator_of_isEquation}; line 163,
\code{affineDegreeEquiv}.\\\bottomrule
\end{longtable}
\normalsize

The list does not contain a packaged canonical degree-one basis, the
precise degree-two kernel, arbitrary infinite arrangements, or the
projective reconstruction theorem. Those are new nodes below. Likewise
the affine-equivalence naturality theorem is not an arbitrary-projection
theorem; the projection wrapper is proved separately on generators.

\Needspace{14\baselineskip}
\section{Coordinates, hyperplanes, and finite index systems}
\label{sec:geometry}

Use $V=k^n$, implemented as \texttt{Fin n} $\to$ \texttt{k}, and
\[
 \cH_n(k)=\{H:\text{\texttt{AffineSubspace k V}}\mid\text{\texttt{IsHyperplane H}}\}.
\]
All incidence statements refer to these affine subspaces. Cohomology
is evaluated on the existing scheme complements, never on their sets
of rational points. Coefficient rings $R$ have their own universe and
no $k$-algebra structure is required.

\node{G01}{HyperplaneGT.coefficientPointEquiv}{existing hyperplane equation API and Mathlib projectivization}
\begin{lemma}\label{lem:g01}
Writing $f_a(x)=a_0+\sum_{i=1}^na_ix_i$, an affine hyperplane is
represented by a tuple $a=(a_0,\ldots,a_n)$ with some $a_i\ne0$ for
$i>0$. Two tuples represent the same hyperplane iff they differ by a
nonzero scalar. Consequently there is a bijection
\[
 \cH_n(k)\simeq\Pj^n(k)\setminus\{o\},\qquad o=[1:0:\cdots:0].
\]
\end{lemma}
\begin{proof}
An affine functional on $k^n$ is its value at zero plus its linear
part, and the latter is the sum of its values on the standard basis
times the coordinates. The linear part is nonzero precisely when one
of those coefficients is nonzero. Existence and uniqueness up to
scalar are \code{exists_isEquation} and \code{existsUnique_smul}.
A projective tuple represents $o$ exactly when all the positive-index
coordinates vanish. The constructions in both directions now have
the asserted formulas, with equality checked by the scalar criterion.
\end{proof}
\emph{Implementation.} Use
\texttt{Projectivization k (Fin (n+1)} $\to$ \texttt{k)}, with the index-zero
coordinate separated by \code{Fin.cases} or \code{Fin.cons}.
\code{Projectivization.mk_eq_mk_iff} and
\code{mk_eq_mk_iff'} in Mathlib's
\code{LinearAlgebra/Projectivization/Basic.lean:112} give the scalar
criterion. Build an \code{Equiv} between the two subtypes and provide
its equation-independence and inverse computation lemmas.

\node{G02}{HyperplaneGT.parallel_iff_normal_proportional}{G01 and the existing codimension/rank formulas}
\begin{lemma}\label{lem:g02}
For distinct hyperplanes $H,K$, either their normal vectors are
proportional and they are disjoint, or their normals are independent
and their intersection has codimension two. For three distinct
hyperplanes $H,K,L$ which are not parallel in pairs, the following
are equivalent: their intersection has codimension two; their affine
equations are linearly dependent; their coefficient points are
collinear. Equal nonempty traces $H\cap K=H\cap L$ are equivalent
to the same dependent-triple condition.
\end{lemma}
\begin{proof}
If the normal rows are proportional, scale them to be equal. Distinct
hyperplanes have different constants, so their equations cannot
vanish simultaneously. If independent, the map $k^n\to k^2$ with
these two rows is onto: otherwise a nonzero linear functional on its
proper image quotient would give a linear relation between the rows.
Solving for the two negative constants gives a point of intersection;
its direction is the joint kernel, of codimension two by rank-nullity.

Translate that point to zero. The linear functionals vanishing on
this joint kernel are exactly the span of the two rows: extend a
basis of the kernel to a basis of $k^n$, or factor through the
surjective two-row map. Thus an affine functional vanishes on the
intersection precisely when it is in the span of the two original
affine equations. For a third distinct hyperplane both coefficients
in that span are nonzero. This proves the triple and trace statements.
Collinearity is membership in the projective span of the two
coefficient vectors, so it is exactly the same condition.
\end{proof}
Define \texttt{Parallel H K} to include $H\ne K$ and disjointness;
define \texttt{DependentTriple H K L} to include pairwise distinctness,
nonempty intersection and codimension two. Give permutation lemmas
once. Do not call an empty-intersection triple dependent.

\node{G03}{HyperplaneGT.vertical_graph_dichotomy}{G01--G02}
\begin{lemma}\label{lem:g03}
For a coordinate subset $J$, a hyperplane is vertical over $J$ iff
its normal coefficients outside $J$ vanish. This is independent of
its equation. If $J$ omits exactly $i$, every nonvertical hyperplane
has a unique equation $x_i-h(z)=0$, $h$ affine linear in $z$.
The projection from this graph to $k^J$ is an affine isomorphism.
If $h_1-h_0$ is a nonzero constant the two graphs are parallel;
if nonconstant they meet, and their common trace projects to
$h_1-h_0=0$.
\end{lemma}
\begin{proof}
Scalar rescaling preserves zero coefficients. Divide the equation
by its nonzero $i$th coefficient and solve for $x_i$. Its inverse
projection inserts $h(z)$ at coordinate $i$. Differences of the
equations give every assertion about intersection; a nonconstant
affine form is onto $k$, so its zero set is nonempty. The graph
equation is unique by comparison of the coefficient at $i$.
\end{proof}
Model coordinate subsets by \texttt{Finset (Fin n)}. Reindex their
coordinate spaces with the subtype of members, avoiding arbitrary
ordered lists. The induction adds a new coordinate to a finite set.

\node{G04}{HyperplaneGT.galoisAction}{G01--G02}
\begin{lemma}\label{lem:g04}
The group $G=(k\simeq_{k_0}k)$ acts on hyperplanes by applying an
automorphism to all equation coefficients. It preserves empty
intersections and their codimensions, parallel pairs, dependent
triples, and verticality over a fixed coordinate subset. Its actions
on equations and projective coefficient points agree.
\end{lemma}
\begin{proof}
Applying $\sigma$ to a scalar multiple gives another scalar multiple,
so G01 makes the operation well defined. Coordinatewise application
$\sigma_V$ satisfies $f_{\sigma a}(\sigma_Vx)=\sigma(f_a(x))$.
It is a bijection, so solutions and emptiness correspond. If
$v_1,\ldots,v_r$ is a basis of a kernel direction, its images span
the transformed kernel: apply $\sigma^{-1}$ to a vector and to
its coefficients. They are independent by the same inverse
argument. Thus dimension and codimension are unchanged. The group
laws are coordinate identities; the remaining conclusions follow
from their definitions.
\end{proof}
Construct a group action, not a postulated natural action on a theory.
Mathlib \code{Projectivization.map} can be used for the semilinear
coordinate map, but the elementary equation argument also suffices.

\node{G05}{HyperplaneGT.finite_orbit}{G04 and finite polynomial root sets}
\begin{lemma}\label{lem:g05}
If $k/k_0$ is algebraic, every hyperplane has a finite $G$-orbit.
Every finite set of hyperplanes therefore has a finite stable
enlargement, and the collection of finite stable arrangements is
cofinal and directed.
\end{lemma}
\begin{proof}
Choose a nonzero normal coefficient and divide it to $1$. This same
coordinate stays $1$ in every conjugate. Each other normalized
coefficient $a$ is algebraic over $k_0$. Its images $\sigma a$ are
roots of the nonzero minimal polynomial of $a$, so form a finite
set. The finite tuple of coefficients ranges in the finite product
of these root sets. Its image in hyperplanes is therefore finite.
For a finite set take the union of its finite orbits; this contains
it and is stable by the group law. Union of two finite stable sets
is a finite stable upper bound.
\end{proof}
Use \code{minpoly.aeval_algHom} in
\code{Mathlib/FieldTheory/Minpoly/Basic.lean:99} and
\code{Polynomial.finite_setOfPred_isRoot} in
\code{Mathlib/Algebra/Polynomial/Roots.lean:154}. Algebraicity implies
integrality over a field and nonzeroness of the minimal polynomial.
No finite normal-closure theorem is needed. The proof itself does
not use perfectness; retain perfectness only in the public
geometric formulation matching the original paper's descended
arrangements.

\node{G06}{HyperplaneGT.CofinalSystem}{G05}
\begin{definition}\label{def:g06}
A cofinal system is a subtype of finite stable arrangements with
the property that each finite set of hyperplanes lies in a member.
Morphisms are inclusions. Markings are the tagged set
$\{(i,s):i<n,s\in S\}\sqcup\{(i,j):i<j<n\}$, with associated
affine functions $x_i-s$ and $x_i-x_j$. For a marking $w$ write
$H_{w,a}=\{w=a\}$ ($a=0,1$). These are distinct hyperplanes and
are individually fixed by $G$.
\end{definition}
\begin{proof}
The functions have nonzero linear part; their levels $0,1$ differ
because $1\ne0$ in $k$. Their coefficients lie in $k_0$, so G04
fixes them. Cofinality gives a member containing both and a common
upper bound for any finitely many selected members. The diagram
is generated by inclusions and a marking arrow $U_A\to T$ when
both labels belong to $A$; impose the evident inclusion/marking
commuting squares. These squares hold as actual morphisms because
all maps are restrictions of the same affine function.
\end{proof}
For the first implementation, use the full type of stable arrangements.
The later cofinal-subsystem theorem is proved by common upper bounds,
not by an assumption that a chosen family contains every union.

\Needspace{14\baselineskip}
\section{The finite degree-one and degree-two interfaces}
\label{sec:finite}

\node{A01}{HyperplaneGT.osDegreeOneEquiv}{existing generator independence and word spanning}
\begin{lemma}\label{lem:a01}
For every finite arrangement $A$ over a field, and any commutative
$R$, the map
\[
 R^{(A)}\longrightarrow\OS_R^1(A),\quad
 (c_H)\longmapsto\sum_Hc_He_H
\]
is a linear equivalence, natural under inclusions. Every element of
the whole \OS{} algebra is a sum of products of these generators
and scalars.
\end{lemma}
\begin{proof}
\code{OSAlgebra.linearIndependent_generator} proves injectivity.
\code{OSAlgebra.span_word_length} at length one says that degree
one is the span of its single-letter words, proving surjectivity.
Both claims hold over the zero ring too. The inclusion formula is
\code{map_generator}; a linear map out of a free module is
determined on its basis. General word spanning gives the last claim.
\end{proof}
Construct the map by \code{Finsupp.linearCombination}; use the
subtype of the degree-one submodule for its target. Package a basis
and its coordinate inverse as corollaries. Avoid \code{finrank},
which would unnecessarily impose field coefficients.

\node{A02}{HyperplaneGT.quadraticKernel_eq_span}{A01, G02, existing relation quotient and grading}
\begin{theorem}\label{thm:a02}
Under A01, the kernel of
$\bigwedge_R^2 R^{(A)}\to\OS_R^2(A)$ is the span of
\[
 e_H\wedge e_K\quad(H\parallel K),\qquad
 (e_H-e_L)\wedge(e_K-e_L)\quad(H,K,L\text{ dependent}).
\]
\end{theorem}
\begin{proof}
Let $\Lambda$ be the exterior algebra on $R^{(A)}$ and let
$J=\Span_R\{arb:r\in\mathcal R,\ a,b\in\Lambda\}$ be its
relation submodule. The actual quotient kernel is $J$, by
\code{OSAlgebra.quotient_eq_zero_iff}. An empty-intersection
relator has degree at least two: the empty list intersects in the
ambient space, and a singleton hyperplane is nonempty. A dependent
boundary has degree at least two: list lengths zero or one are
impossible, and distinct intersecting hyperplanes have codimension
two by G02, ruling out a dependent list of length two. The existing
private lemma \code{dependent_length_ge_three} in
\code{Presentation/Generators.lean} contains the latter proof;
expose a suitable public lemma or repeat its short list-length argument.

For a finite sum of contexts $\sum a_jr_jb_j$, decompose each
$a_j,b_j$ in the nonnegative exterior grading. Multiplication adds
degrees. Projection to degree two has a contribution only if the
relator has degree two and both contexts have degree zero. Degree
zero is the scalar copy of $R$ in the exterior algebra. Hence
$J\cap\Lambda^2$ is exactly the scalar span of degree-two relators.
They are length-two empty words and length-three dependent
boundaries. By G02 these are precisely the displayed configurations.
Expand
\[
 (e_H-e_L)(e_K-e_L)=e_He_K-e_He_L+e_Ke_L
\]
using $e_L^2=0$ and $e_Le_K=-e_Ke_L$. This is the ordered triple
boundary. The second exterior power identifies with $\Lambda^2$;
thus the kernel is exactly as claimed. No assumption about higher
relations being generated by degree-two relations has been made.
\end{proof}
\emph{Implementation detail.} Prove a generic lemma for a
nonnegatively graded algebra with relators of degree $\geq d$:
its relation submodule in degree $d$ is spanned by degree-$d$
relators with degree-zero left/right contexts. Prove it by
\code{Submodule.span_induction} on contexts and homogeneous
projection, using the finite support of each decomposition. Apply
it at $d=2$. Use Mathlib
\code{exteriorPower.alternatingMapLinearEquiv} for alternating maps.
In this Mathlib revision, \code{ExteriorAlgebra.exteriorPower} is
definitionally the submodule
\texttt{(LinearMap.range (ExteriorAlgebra.}$\iota$\texttt{ R)) \^{} n}
of the exterior
algebra (\code{LinearAlgebra/ExteriorAlgebra/Basic.lean:79}), and
\code{ExteriorAlgebra.gradedAlgebra} uses these pieces
(\code{ExteriorAlgebra/Grading.lean:82}). Thus the second exterior
power already is the required degree-two submodule. The coefficient
formula \code{exteriorPower.}$\iota$\code{Multi_apply_coe} in
\code{ExteriorPower/Basic.lean:60} identifies its wedge with the
same exterior word; no new comparison isomorphism is needed.

\node{A03}{HyperplaneGT.finiteQuadraticDiagram}{A01--A02 and G06}
\begin{definition}\label{def:a03}
Set $M_A=R^{(A)}$ and let $Q_A$ be the submodule in A02. For the
extra object $T$, set $M_T=R^{\{0,1\}}$ and $Q_T=\bigwedge^2M_T$.
Inclusions send a basis label to the same label. Marking pullback
sends the two ordered basis vectors to $e_{H_{w,0}},e_{H_{w,1}}$.
The latter sends $Q_T$ into $Q_A$ because its two labels are parallel.
A quadratic diagram automorphism consists of linear equivalences
$g_A,g_T$ commuting with these maps and satisfying
$(\bigwedge^2g_A)Q_A=Q_A$.
\end{definition}
Use the subgroup of the product of the indicated \code{LinearEquiv}
groups cut out by these equations. Composition and inverse
preserve the equations: for inverse naturality multiply the
naturality square by the two inverse maps. Equality of submodule
images is necessary; mere containment is insufficient for the
inverse. This construction uses finite combinatorial modules and
does not assert any cohomological conclusion as an input.

\node{A04}{HyperplaneGT.degreeOneDiagramEquiv}{A01--A03 and the existing diagonal structure theorem}
\begin{theorem}\label{thm:a04}
For a supplied C1--C5 theory, its degree-one and cup-kernel diagram
is isomorphic to A03. The isomorphism sends $e_H$ to $[f_H]$ and
the ordered basis of $M_T$ to $[t]$, $[t-1]$.
Moreover the canonical diagonal algebra equivalences form an
isomorphism of the full graded diagrams.
\end{theorem}
\begin{proof}
Compose A01 with
\code{Geometric.OrdinaryTheory.degreeEquiv} at degree one. Since the full
\code{diagonalEquiv} preserves products and is bijective, its
degree-two restriction identifies the two product kernels. For
$T$, both points are parallel and A02 says that its whole second
exterior power is the kernel; A01 gives its stated basis.

Inclusion naturality follows by the existing generator and
\code{affineEquiv_map} formulas. For a marking $w$, the coordinate
map to $\A^1$ restricts to a scheme map into $T$ because $w$ and
$w-1$ are units on the source complement. This is the localization
universal property in \code{SmoothAlg.complementMap}, transported
by the existing geometric affine comparison. Logarithmic naturality
gives $w^*[t-a]=[w-a]$. It proves the square on both basis vectors
and therefore everywhere. Identity and composition follow on
these same generators. Conjugating families by these linear
equivalences yields a multiplicative equivalence of automorphism
groups and carries scalar families to scalar families.
The same generator comparison proves naturality of the full
diagonal algebra equivalences: inclusions agree on every
generator, and each marking agrees on both generators of the
algebra of $T$. Since these algebras are generated in degree one,
the comparison is an isomorphism of the full graded diagrams too.
\end{proof}

\Needspace{14\baselineskip}
\section{The universal free module and its dual}
\label{sec:universal}

\node{U01}{HyperplaneGT.universalDegreeOne}{G06 and A03}
\begin{definition}\label{def:u01}
Define $M=R^{(\cH_n(k))}$, the finitely supported functions on all
hyperplanes, with $e_H$ the standard basis. Include each $M_A$ by
extension by zero. Its image is exactly the vectors supported on
$A$. Define $Q\subset\bigwedge^2M$ as the span of all parallel-pair
and dependent-triple relators.
\end{definition}
\begin{lemma}\label{lem:u01}
These inclusions exhibit $M$ as the direct limit of the finite
degree-one diagram. Every finite sum of generators of $Q$ is the
image of an element of $Q_A$ for some $A\in\mathscr F$, and every
image of $Q_A$ is contained in $Q$.
\end{lemma}
\begin{proof}
A vector has finite support, contained in some $A$ by cofinality.
Two representatives have the same coefficients after inclusion in
a common upper bound, proving the direct-limit description. More
explicitly, a compatible family of linear maps from $M_A$ to any
module $N$ gives a function on the labels: choose a stage containing
$H$ and evaluate its map at $e_H$. Common upper bounds make this
independent of the stage. Extend by the finite sum
$\sum_Hc_H f(e_H)$. This proves the universal property and uniqueness.
For a sum of relators take a stage containing the union of the two
or three labels in every summand. Each is then one of A02's
relators there. Conversely those relators are among the global
ones, so their image and its linear span lie in $Q$.
\end{proof}
No injectivity of degree-two restriction maps is required for this
argument. Define the global relation module by this span, rather
than invoking an unproved assertion about a cohomological colimit.

\node{U02}{HyperplaneGT.DiagramAut.universalEquiv}{U01 and A03}
\begin{lemma}\label{lem:u02}
A quadratic-diagram automorphism induces a linear equivalence of
$M$, preserving $Q$. It preserves each finite-stage submodule, and
the marked maps from $R^2$ have the common equivariant source
$(R^2,g_T)$. The construction is a group homomorphism. The induced
universal map together with $g_T$ determines the original family.
\end{lemma}
\begin{proof}
Use U01 on the maps $M_A\xrightarrow{g_A}M_A\to M$. Naturality
gives compatibility. The family of inverses is compatible too;
its universal map is an inverse because both composites fix each
finite-stage vector. For a global relation choose a finite stage
as in U01; its image under $g$ lies in the image of $Q_A$, hence
in $Q$. Apply the same argument to $g^{-1}$ for equality. On the
two marked basis vectors the equivariance is the finite-stage
naturality equation. Group multiplication and injectivity are
checked on vectors in finite stages.
\end{proof}
The auxiliary universal construction must not be declared as the
group of all automorphisms of $M$ preserving $Q$: it also retains
finite-stage invariance and the common marked source. Those data
encode the arithmetic diagram.

\node{U03}{HyperplaneGT.dualCoordinates}{U01 and Mathlib free-module duality}
\begin{lemma}\label{lem:u03}
There is an explicit linear equivalence
\[
 P:=\Hom_R(M,R)\simeq\prod_{H\in\cH_n(k)}R,
 \qquad\lambda\longmapsto(\lambda(e_H))_H.
\]
Its inverse sends a function $a$ to
$m\mapsto\sum_{H\in\operatorname{supp}m}m_Ha_H$.
Evaluation separates points on both sides. A linear equivalence
$g$ of $M$ acts on $P$ by $g\lambda=\lambda\circ g^{-1}$, and
$(g\lambda)(gm)=\lambda(m)$.
\end{lemma}
\begin{proof}
Finite support makes the formula well defined and linear. Its value
on $e_H$ is $a_H$, and expansion of any $m$ in the basis proves
the other composite identity. Distinct vectors are separated by
some coordinate functional $\delta_H$; distinct functionals differ
on some $e_H$. The inverse in the displayed group action makes it
a left action: $(gh)\lambda=\lambda\circ h^{-1}\circ g^{-1}$,
which equals $g(h\lambda)$. The evaluation identity is substitution.
\end{proof}
Use \texttt{Finsupp.llift R R R X}, where $X=\cH_n(k)$, in
\code{Mathlib/LinearAlgebra/Finsupp/LSum.lean:157}, with
\code{llift_symm_apply} at line 171, or the same formulas directly.
\code{LinearEquiv.dualMap} is available in
\code{LinearAlgebra/Dual/Defs.lean:141}; verify its direction and
compose with the inverse as prescribed. No topology and no
double-dual equivalence are used.

\node{U04}{HyperplaneGT.Commutes}{U03 and U01}
\begin{definition}\label{def:u04}
For $\lambda,\mu\in P$ define the alternating form
$d_{\lambda,\mu}(a,b)=\lambda(a)\mu(b)-\lambda(b)\mu(a)$.
Its exterior lift is a map $\bigwedge^2M\to R$. Define
\[
 \lambda\sim\mu\quad\Longleftrightarrow\quad
 Q\subseteq\ker d_{\lambda,\mu},\qquad
 C(\lambda)=\{\mu\in P:\lambda\sim\mu\}.
\]
\end{definition}
\begin{lemma}\label{lem:u04}
$C(\lambda)$ is a submodule. The commuting relation is equivalent
to the following equations on all the indicated configurations:
\begin{align*}
 \lambda_H\mu_K&=\lambda_K\mu_H&& (H\parallel K),\\
 (\lambda_H-\lambda_L)(\mu_K-\mu_L)
 &=(\lambda_K-\lambda_L)(\mu_H-\mu_L)
 &&(H,K,L\text{ dependent}).
\end{align*}
It is preserved in both directions by the paired action in U02--U03.
\end{lemma}
\begin{proof}
Alternation is $\lambda(a)\mu(a)-\lambda(a)\mu(a)=0$, so the
exterior lift exists even in characteristic two. The formula is
linear in $\mu$, proving the submodule assertion. Vanishing on a
span is equivalent to vanishing on its generators, and evaluation
on the two types of relator is precisely the display. Finally
\[
 d_{g\lambda,g\mu}(r)=d_{\lambda,\mu}((\bigwedge^2g^{-1})r).
\]
It holds on simple wedges by evaluation and hence everywhere by
their spanning property. Preservation of $Q$ by $g$ and its inverse
proves the equivalence.
\end{proof}
Use \code{exteriorPower.alternatingMapLinearEquiv} or
\code{exteriorPower.pairingDual} and
\code{alternatingMapToDual_apply_}$\iota$\code{Multi} in
\code{Mathlib/LinearAlgebra/ExteriorPower/Pairing.lean:58--64}.
The latter is a determinant pairing, not a claim that the dual of
an infinite exterior power is a particular tensor product.

\node{U05}{HyperplaneGT.primitive_iff_unit_coordinate}{U03 and local-ring arithmetic}
\begin{lemma}\label{lem:u05}
For a nonzero local ring $R$, a functional $\lambda:M\to R$ is
surjective iff some $\lambda_H$ is a unit. Surjectivity is invariant
under precomposition by any linear equivalence of $M$.
\end{lemma}
\begin{proof}
A preimage of $1$ gives a finite sum
$1=\sum_Hm_H\lambda_H$. In a local ring the nonunits form the
maximal ideal, so at least one summand is a unit, and then its
factor $\lambda_H$ is a unit. Conversely, if $\lambda_H=u$ is a
unit, $u^{-1}e_H$ maps to $1$ and $ru^{-1}e_H$ maps to any $r$.
Precomposition with a bijection has the same image.
\end{proof}
Use \code{IsLocalRing.exists_of_isUnit_sum} in
\code{Mathlib/RingTheory/LocalRing/Basic.lean:58} and the fact that
a factor of a unit product in a commutative ring is a unit. Define
\texttt{Primitive} $\lambda$ to mean \texttt{Function.Surjective} $\lambda$. Do not
replace it by $\lambda\notin\mathfrak mP$, which has different
behavior for non-finitely-generated maximal ideals in an infinite
product. This is the only use of localness in the core reconstruction.

\Needspace{14\baselineskip}
\section{Trace modules, centralizers, and local detection}
\label{sec:local}

\node{L01}{HyperplaneGT.traceMap}{G02--G03 and U03}
\begin{lemma}\label{lem:l01}
For $H\in\cH_n(k)$, each distinct $K$ has either empty trace on
$H$ or a hyperplane trace. Every hyperplane of $H$ is such a trace.
Define $U_H\subset M$ by coefficient zero at $H$, and define
\[
 s_H:U_H\longrightarrow R^{(\cH(H))},\qquad
 e_K\longmapsto
 \begin{cases}0&K\parallel H,\\ e_{H\cap K}&\text{otherwise}.
 \end{cases}
\]
Then $s_H$ is onto and its kernel $U_H^1$ is spanned by
\[
 e_K\quad(K\parallel H),\qquad
 e_K-e_L\quad(H,K,L\text{ a dependent triple}).
\]
\end{lemma}
\begin{proof}
The trace dichotomy is G02. Choose a point in $H$ and a vector
transverse to its direction. The map that subtracts the transverse
coordinate is an affine retraction $V\to H$. Pulling a nonconstant
affine equation of a hyperplane of $H$ back along that retraction
gives an equation of a hyperplane $K\ne H$ with exactly that trace.
This proves surjectivity on labels and therefore on the free module.

For the kernel, let $m=\sum_{K\ne H}c_Ke_K$ have finite support.
Separate empty-trace terms. In each remaining trace fiber $J$
choose one representative $K_J$ from the support, if the fiber is
nonempty. Then
\[
 \sum_{K:\,H\cap K=J}c_Ke_K
 =\sum_{K:\,H\cap K=J}c_K(e_K-e_{K_J})
   +\Bigl(\sum_{K:\,H\cap K=J}c_K\Bigr)e_{K_J}.
\]
If $s_H(m)=0$, every final coefficient is zero because these are
the coefficients of distinct basis vectors $e_J$. The differences
with $K=K_J$ are zero; all other differences are from dependent
triples by G02. This proves generation of the kernel. Each stated
generator visibly maps to zero, proving the reverse inclusion.
\end{proof}
The intrinsic affine space $H$ can be represented by the existing
chart on \code{H.direction}; do not redefine all arrangement
machinery for torsors. The statement and subsequent maps are
independent of that chart because they only mention trace equality.
Factor the finite-fiber kernel argument as a general free-module
lemma for a map of label sets with an optional discarded value.

\node{L02}{HyperplaneGT.inertiaDecomposition}{L01 and U03--U04}
\begin{definition}\label{def:l02}
Put $I_H=R\delta_H$ and let $D_H$ consist of functions whose
value is zero on hyperplanes parallel to $H$ and equal on any two
distinct hyperplanes with the same nonempty trace on $H$.
\end{definition}
\begin{lemma}\label{lem:l02}
With annihilators taken under the evaluation pairing,
\[
 I_H=\Ann(U_H),\quad D_H=\Ann(U_H^1),\quad
 \Ann(I_H)=U_H,\quad\Ann(D_H)=U_H^1.
\]
There is an explicit equivalence
$D_H\simeq R\times R^{\cH(H)}$, recording the value at $H$ and
the value on each trace. Its second projection has kernel $I_H$.
For a unit $u\in R^\times$, $C(u\delta_H)=D_H$.
\end{lemma}
\begin{proof}
The first two annihilator statements follow by evaluation on L01's
generators. For the product equivalence the inverse assigns value
$a$ at $H$, zero on empty traces, and $b_J$ on trace $J$. The
surjectivity on trace labels in L01 makes the two formulas inverse.
The kernel of the second projection is exactly the arbitrary value
at $H$ alone, namely $I_H$.

For $\Ann(D_H)=U_H^1$, first evaluate at $\delta_H\in D_H$ to
force the coefficient at $H$ to be zero. If $s_H(m)\ne0$, choose
a trace with nonzero coefficient and the function in $D_H$ that
is its indicator and has value zero at $H$. Its value on $m$ is
that coefficient, a contradiction. Thus $s_H(m)=0$. The converse
follows directly from the first two identities. Similarly
$\Ann(I_H)=U_H$ is evaluation at $\delta_H$.

For the centralizer formula, U04's tests with $H$ absent are zero.
For a parallel pair containing $H$ they say $u\mu_K=0$, hence
$\mu_K=0$. For a dependent triple containing $H$ they say
$u(\mu_K-\mu_L)=0$, hence $\mu_K=\mu_L$. The converse is given
by the same tests. This uses a unit cancellation, not domain
cancellation. No reflexivity or exactness theorem for algebraic
duals of arbitrary modules has been invoked.
\end{proof}

\node{L03}{HyperplaneGT.graph_decomposition_projection}{L02 and G03}
\begin{lemma}\label{lem:l03}
Let $z$ omit coordinate $i$, and let $M_z$ be the free span in
$M$ of vertical hyperplanes over $z$. Set $P_z=\Hom_R(M_z,R)$
and $p_z:P\to P_z$ to be restriction. If $H$ is a graph over $z$,
then $p_z(D_H)=P_z$ and the restriction has kernel $I_H$.
Moreover $p_z(\delta_H)=0$.
\end{lemma}
\begin{proof}
The graph is affine isomorphic to $k^z$ by G03. Each hyperplane
of $k^z$ pulls back to one vertical hyperplane of $V$, whose trace
on $H$ is that hyperplane under the isomorphism. Conversely these
are all hyperplanes of $H$, with no repetitions and no empty
traces. Thus the product description in L02 identifies $p_z$
with the projection to its second factor; its first factor is
$I_H$. The graph label is not vertical, proving the last assertion.
\end{proof}

\node{L04}{HyperplaneGT.local_transfer}{G03, U04 and surjectivity of a centralizer projection}
\begin{lemma}\label{lem:l04}
Suppose $p_z\lambda=0$, $p_z:C(\lambda)\to P_z$ is onto, and
$H_0$ is a graph with $\lambda_{H_0}$ a unit. For a graph $H_1$
distinct and parallel to $H_0$ and a graph $H_2\ne H_0$ meeting
$H_0$, one has $\lambda_{H_1}=\lambda_{H_2}$.
\end{lemma}
\begin{proof}
Use the unique graph equations to write
$H_0:x_i=h_0(z)$, $H_1:x_i=h_0(z)+c$, and
$H_2:x_i=h_2(z)$, with $c\ne0$ and $h_2-h_0$ nonconstant.
Let $V_1:h_2-h_0=0$ and $V_2:h_2-h_0=c$. They are distinct
vertical hyperplanes. The triples $(H_0,H_2,V_1)$ and
$(H_1,H_2,V_2)$ have nonempty codimension-two intersection:
solve the nonconstant affine equation in $z$, then insert the
graph coordinate. They are pairwise distinct and hence dependent.
Choose a function on vertical labels taking values $0$ at $V_1$
and $1$ at $V_2$; set it zero elsewhere. By surjectivity it lifts
to $\mu\in C(\lambda)$. Put $a_j=\lambda_{H_j}$,
$b_j=\mu_{H_j}$. Since $\lambda_{V_1}=\lambda_{V_2}=0$,
U04 gives
\[
 a_0b_1=a_1b_0,\qquad a_0b_2=a_2b_0,\qquad
 a_1(b_2-1)=a_2(b_1-1).
\]
Writing $b_1=a_0^{-1}a_1b_0$, $b_2=a_0^{-1}a_2b_0$ in the
last equality cancels their common product and yields $a_1=a_2$.
The inverse denotes the inverse of the specified unit, not a
field inverse on $R$.
\end{proof}
Factor the last calculation into a commutative-ring lemma with a
unit parameter and the three equations as hypotheses. This makes
the local-ring and geometry proof independent of ring-tactic
normalization details.

\node{L05}{HyperplaneGT.localDetection}{U05, L02--L04 and G03}
\begin{theorem}\label{thm:l05}
Assume $n\geq2$ and $R$ is nonzero local. A functional $\lambda$
is $u\delta_H$ for a unit $u$ and a graph $H$ over $z$ iff:
\begin{enumerate}
\item $\lambda$ is primitive;
\item $p_z\lambda=0$;
\item there is a graph $H'$ over $z$ with $\lambda_{H'}=0$;
\item $p_z:C(\lambda)\to P_z$ is surjective.
\end{enumerate}
The graph $H$ is unique, and then $C(\lambda)=D_H$.
\end{theorem}
\begin{proof}
For $u\delta_H$, primitivity follows by evaluation, vertical
vanishing and centralizer surjectivity follow from L02--L03,
and any different graph supplies a zero coordinate.

Conversely U05 gives a unit coordinate $H_0$. Vertical vanishing
and nonzero coefficients make it a graph. Distinct parallel graphs
exist by replacing $h_0$ with $h_0+1$. Distinct intersecting graphs
exist by replacing $h_0$ with $h_0+z_j$ for any remaining
coordinate $j$; one exists because $n\geq2$.
The zero-coordinate graph $H'$ is distinct from $H_0$. If it is
parallel, retain it. Otherwise L04 compares it with a parallel
graph, giving a parallel graph $H_1$ of value zero. Now L04 makes
every graph distinct from and meeting $H_0$ have value zero.
Any other graph parallel to $H_0$ compares with a chosen
intersecting graph and also has value zero. Vertical labels are
already zero. Thus $\lambda=\lambda_{H_0}\delta_{H_0}$.
If two such unit-supported functions agree, evaluation at their
support labels forces equality of the labels and units, proving
uniqueness. L02 gives the centralizer assertion.
\end{proof}
Only surjectivity of the centralizer projection was used. The
stronger quotient-isomorphism condition in the original paper is
a consequence by L03, not an additional assumption here.

\Needspace{14\baselineskip}
\section{Recognizing the omitted coordinate projections}

\node{P01}{HyperplaneGT.coordinateSpan_recognition}{G02--G03, U04 and L02}
For a coordinate subset $z\subset\{1,\ldots,n\}$ let $M_z$ be
the span of all hyperplanes whose equations use only $z$, and let
$P_z$ be the product indexed by these hyperplanes. Restriction
$p_z:P\to P_z$ is surjective, by extension by zero, and
$\ker p_z=\Ann(M_z)$. For $w\notin z$ put
$L_w=Re_{w=0}+Re_{w=1}$.
\begin{lemma}\label{lem:p01}
An element $a\in M$ belongs to $M_{z\cup\{w\}}$ iff
\[
 d_{\lambda,\mu}(a\wedge b)=0
 \quad\text{for every }b\in L_w\text{ and every commuting }
 \lambda,\mu\in\Ann(M_z).
\]
\end{lemma}
\begin{proof}
It suffices in the forward direction to take a basis vector $e_H$
and $b=e_{w=t}$, $t=0,1$. If $H$ is vertical over $z$, both
functionals vanish on $e_H$. Otherwise its equation can be normalized
as $w+L(z)+c=0$. If $L=0$, $H$ equals or is parallel to $w=t$,
and the determinant is zero by alternation or U04. If $L\ne0$,
the three distinct hyperplanes
\[
 w+L(z)+c=0,\qquad w=t,\qquad L(z)+c+t=0
\]
have a common codimension-two intersection. The last is vertical
over $z$. The triple equation of U04, with its two functional
values zero there, is the required determinant equation.

Conversely suppose $H$ is not vertical over $z\cup\{w\}$.
It is neither equal nor parallel to $w=0$: either would supply
an equation in $w$ alone. Thus $J=H\cap\{w=0\}$ is a hyperplane
of $H$. This trace is different from the nonempty trace of every
hyperplane $V$ vertical over $z$. Here is the required equation
argument. If the traces agree, the restrictions to $H$ of $w$ and
an equation $v$ of $V$ are nonzero affine forms with the same zero
hyperplane, hence are scalar multiples by G01 applied on $H$.
For some $c\ne0$, $w-cv$ vanishes on $H$. It is nonzero, since
$v$ does not involve $w$. An affine form vanishing on $H$ is a
multiple of its equation: translate to a point of $H$, choose a
vector transverse to its direction, and evaluate on this direct
sum. Therefore $H$ has equation $w-cv$, a contradiction.

Take $\lambda=\delta_H$ and the element $\mu\in D_H$ whose
coefficient at $H$ is zero, whose trace value at $J$ is one, and
whose other trace values are zero. Both annihilate $M_z$, by the
preceding trace assertion; they commute by L02. Moreover
$\lambda(e_{w=0})=0$ and $\mu(e_{w=0})=1$. The displayed test
with $b=e_{w=0}$ consequently reads $a_H=0$. Apply this to every
label outside the coordinate span. A finitely supported vector
with all those coefficients zero belongs to $M_{z\cup\{w\}}$,
by its basis expansion.
\end{proof}
\emph{Formal details.} Prove the support characterization of $M_z$
once using the Finsupp basis. For the equation argument on $H$, an
affine coordinate equivalence $H\simeq\A^{n-1}$ exists by extending
a basis of its direction subspace and translating by a point.
Alternatively prove the affine-form restriction statement directly
by the transverse-vector argument above; no chosen coordinate
equivalence then escapes the lemma.

\node{P02}{HyperplaneGT.coordinateSpan_invariant}{P01, U02--U04 and A03}
\begin{proposition}\label{prop:p02}
Every diagram automorphism preserves $M_z$ for every coordinate
subset $z$. Its dual induces an equivalence of $P_z$ making $p_z$
equivariant.
\end{proposition}
\begin{proof}
Induct on the finite cardinality of $z$. The empty span is zero.
Choose $w\in z$ at the induction step. The marked map $w:T$
identifies its degree-one module with $L_w$, and naturality makes
$L_w$ invariant. The induction hypothesis makes $M_{z\setminus\{w\}}$
invariant. Evaluation and commuting are invariant by U03--U04,
so P01 gives forward containment for $M_z$. Apply the same
argument to the inverse to obtain equality.

Since $\ker p_z=\Ann(M_z)$ is invariant, define the induced map on
$P_z$ by choosing an extension, applying the dual, and restricting.
Two extensions differ by the kernel, so the result is independent
of the choice. The inverse construction gives an equivalence;
the construction is linear and satisfies the equivariance equation.
\end{proof}
This quotient argument involves only the explicitly surjective
restriction map. It uses no exactness of dualization.

\Needspace{14\baselineskip}
\section{Hyperplane permutation and the projective incidence structure}

Throughout I01--I04, C04 and S01--S04, the coefficient ring $R$
is nonzero and local. The intervening pure geometric lemmas
C01--C03 only require a field $k$ and $n\ge2$. S05 separately
extends the reconstruction theorem to integral domains.

\node{I01}{HyperplaneGT.DiagramAut.hyperplanePerm}{L05 and P02}
\begin{proposition}\label{prop:i01}
Each diagram automorphism $g$ has a unique permutation $\pi_g$ of
$\cH$ such that on the dual $gI_H=I_{\pi_gH}$ and
$gD_H=D_{\pi_gH}$. On $M$ it satisfies
$g(Re_H)=Re_{\pi_gH}$. These permutations respect composition.
\end{proposition}
\begin{proof}
Choose a nonzero normal coefficient of $H$, say its $i$th, so $H$
is a graph over the complementary coordinates $z$. The functional
$\lambda=g\delta_H$ is primitive because $g$ is an equivalence.
The vanishing on $M_z$ and surjectivity of its centralizer
projection follow from L03, P02, and preservation of commuting.
To obtain the zero-coordinate graph required in L05, first suppose
$H\notin\{x_i=0,x_i=1\}$. Then $\delta_H$ annihilates the
invariant plane $L_{x_i}$; so does $\lambda$, supplying two graph
coordinates of value zero. If $H$ is one of these two hyperplanes,
choose $j\ne i$. The common marking $x_r-x_s$, where
$\{r,s\}=\{i,j\}$ and $r<s$, has two graph hyperplanes distinct
from $H$. Their invariant plane is annihilated by $\delta_H$
and hence by $\lambda$. This supplies the required zero coordinate.
L05 gives a unique graph $K$ and a unit $u$ with
$\lambda=u\delta_K$. The construction is independent of $i$
because distinct coordinate inertia lines are distinct over a
nonzero ring. Applying the construction to $g^{-1}$ proves it is
a permutation. L02 and equivariance of centralizers give the
assertion about $D_H$.

The annihilator identity $U_H=\Ann(I_H)$ is equivariant, and
\[
 Re_H=\bigcap_{K\ne H}U_K.
\]
Indeed membership says that every coefficient other than $H$ is
zero; a Finsupp vector then equals $a_He_H$. This proves the
generator-line assertion. Equality of the permutations of a
composition follows from uniqueness on inertia lines.
\end{proof}
The action on $M$ is $g$ and the action on $P$ is the
contragredient $\lambda\mapsto\lambda\circ g^{-1}$. Keeping this
convention throughout makes both permutations $\pi_g$, not inverses.

\node{I02}{HyperplaneGT.DiagramAut.preserves_incidence}{I01 and L01--L02}
\begin{lemma}\label{lem:i02}
The permutation $\pi_g$ preserves parallel pairs and dependent
triples, in both directions.
\end{lemma}
\begin{proof}
The equality $U_H^1=\Ann(D_H)$ from L02 implies
$gU_H^1=U_{\pi_gH}^1$. For distinct $H,K$,
\[
 H\parallel K\quad\Longleftrightarrow\quad Re_K\subset U_H^1.
\]
If the trace is empty its image is zero; otherwise its image is a
basis vector, nonzero because $R$ is nonzero. For distinct $H,K,L$,
dependence is equivalent to the images of $Re_K,Re_L$ in
$U_H/U_H^1$ being equal nonzero lines. By L01 this quotient is the
free module on traces. Two coordinate lines in a free module over
a nonzero ring agree only if their labels agree: evaluate a basis
vector at its own coordinate. G02 identifies equal nonempty traces
with dependence. These characterizations are preserved by the
equivalences and their inverses.
\end{proof}

\node{I03}{HyperplaneGT.DiagramAut.projectiveCollineation}{G01--G02 and I02}
Identify $\cH$ with $\Pj^n(k)\setminus\{o\}$ via coefficient
points, where $o=[1:0:\cdots:0]$. Extend $\pi_g$ to fix $o$.
\begin{lemma}\label{lem:i03}
This extended permutation maps each projective line onto a
projective line, and its inverse does the same.
\end{lemma}
\begin{proof}
A line through $o$ has vector space span of $e_0$ and an affine
form with nonzero normal part. Its other points have proportional
normal parts and arbitrary constant, exactly a parallel pencil.
Conversely such a pencil is this line minus $o$. I02 applies.

For a line not through $o$, take independent representatives $f,g$.
Their normal parts are independent: a nontrivial dependence of
normal parts would give a nonzero constant in their span (zero
would contradict independence of $f,g$). The associated linear
map $k^n\to k^2$ is surjective: if its image were proper, a
nonzero linear functional annihilating that image, obtained by
extending a basis, would give a dependence of the two rows.
Thus the two affine zero sets meet in codimension two. The other
line points have representatives $af+bg$ with $ab\ne0$, and form
dependent triples with them. Conversely an affine form vanishing
on their common intersection is in their span: translate to a
common point and factor the linear form through the surjective
map $k^n\to k^2$ by its kernel. I02 consequently preserves exactly
the other points of this line. Applying the same argument to the
inverse proves equality of line images.
\end{proof}
\emph{Mathlib interface.} Use \code{Projectivization} of
$k^{n+1}$ and projectivizations of two-dimensional submodules.
The existing \code{Projectivization.IsCollinear} and
\code{Projectivization.line_unique} are in
\code{Mathlib/LinearAlgebra/Projectivization/Collinear.lean:41,112}.
The coefficient model also permits proving this lemma with spans
first, then packaging the projectivization statement.

\node{I04}{HyperplaneGT.marked_hyperplanes_fixed}{I01--I02 and A03}
\begin{lemma}\label{lem:i04}
Each marked hyperplane $w=0$ and $w=1$ is individually fixed.
\end{lemma}
\begin{proof}
Invariance of the free coordinate plane $L_w$ and I01 imply that
its two labels are permuted. In its ordered basis the matrix of
$g$ is the same matrix $g_T$ for every $w$, by naturality and
injectivity of the marked map. A monomial invertible two-by-two
matrix over a nonzero ring has a unique support permutation:
its selected entries are units and hence nonzero. Thus either
every marked pair is fixed or every pair is exchanged. Choose
$i<j$. The triple $x_i=0,x_j=0,x_i-x_j=0$ is dependent. A common
exchange would send it to $x_i=1,x_j=1,x_i-x_j=1$, whose
intersection is empty since $1-1=0\ne1$. This contradicts I02.
The argument concerns permutations of labels, so it also works
when $R$ or $k$ has characteristic two.
\end{proof}

\Needspace{14\baselineskip}
\section{A coordinate proof of field reconstruction}

\node{C01}{HyperplaneGT.collineation_preserves_span}{an arbitrary line-preserving bijection and finite-dimensional linear algebra}
\begin{lemma}\label{lem:c01}
A projective bijection preserving lines in both directions
preserves the projective span of every finite set.
\end{lemma}
\begin{proof}
The projective span is the smallest set containing the given
points and closed under joining two distinct points by their
entire projective line. To verify the assertion, take a nonzero
linear combination of representatives and induct on the number
of summands. Discard zero summands. If the partial sum is zero,
the last summand alone represents the result. If it is nonzero
and proportional to the last summand, a nonzero total represents
the same already constructed point. Otherwise the two partial
points are distinct, and the total lies on their line. This
proves the characterization. Apply it to the bijection and its
inverse to obtain equality of spans.
\end{proof}
The existing span objects are
\code{Projectivization.Subspace.span} in
\code{Mathlib/LinearAlgebra/Projectivization/Subspace.lean:85}
and \code{Submodule.projectivization} at line 285. The finite
sum proof above avoids any assumed fundamental theorem of
projective geometry.

\node{C02}{HyperplaneGT.collineation_coordinate_functions}{C01}
\begin{lemma}\label{lem:c02}
Suppose $F$ is a projective bijection preserving lines in both
directions and fixes $[e_0],[e_i]$ for
$1\le i\le n$. Then on the chart $a_0=1$ it has the form
\[
 F(a_1,\ldots,a_n)=(f_1(a_1),\ldots,f_n(a_n))
\]
with bijections $f_i:k\to k$ and $f_i(0)=0$.
\end{lemma}
\begin{proof}
C01 makes the hyperplane $a_0=0$, spanned by the fixed points
$[e_i]$, invariant. Thus the chart is invariant. Its origin
$[e_0]$ is fixed. An affine line parallel to axis $i$ is a
projective line with unique point at infinity $[e_i]$, so its
image has the same direction. Consequently, changing coordinate
$i$ of a point changes no other coordinate of its image.
For each output coordinate $j$, change the other input
coordinates successively to zero. Its value is unchanged and
therefore depends only on input coordinate $j$. Define $f_j$ by
restriction to that axis. C01 also preserves the axis itself,
and the inverse restricts to its inverse, proving bijectivity.
The fixed origin gives the asserted zero values.
\end{proof}

\node{C03}{HyperplaneGT.coordinateFieldEquiv}{C02, the marked frame and $n\ge2$}
\begin{theorem}\label{thm:c03}
If $F$ additionally fixes $[e_i-e_0]$ for every $i$ and
$[e_i-e_j]$ for $i<j$, there is a unique field automorphism
$\sigma$ such that $F([a])=[\sigma(a)]$ coordinatewise.
\end{theorem}
\begin{proof}
The first additional fixed points are $-e_i$ in the affine
chart, so $f_i(-1)=-1$. In the $(i,j)$ coordinate plane the
affine lines $a_i+a_j=c$ have point at infinity $[e_i-e_j]$;
the plane is invariant by C01 and their direction is fixed.
For $c\ne0$, the two distinct points $ce_i,ce_j$ on such a
line have images $f_i(c)e_i,f_j(c)e_j$, and constancy of the
coordinate sum along the image line gives $f_i(c)=f_j(c)$.
For $c=0$ both values are zero by C02. Thus all $f_i$ are a
single bijection $\sigma$. Every point $ae_i+be_j$ on the
line with $c=a+b$ has image coordinate sum $\sigma(a+b)$,
so $\sigma(a)+\sigma(b)=\sigma(a+b)$. This also holds when
$c=0$, since that line maps to the line of the same direction
through the fixed origin. Additivity and $\sigma(-1)=-1$
give $\sigma(1)=1$.

The line through $0$ and $e_i+be_j$ has points
$ae_i+abe_j$. Its image passes through $0$ and
$e_i+\sigma(b)e_j$, so comparison of its two coordinates gives
$\sigma(ab)=\sigma(a)\sigma(b)$. Package these identities and
bijectivity as a ring equivalence of $k$ with itself.

For $v\ne0$, the projective line through $[1:0]$ and $[1:v]$
has unique point at infinity $[0:v]$. The image line has
unique point at infinity $[0:\sigma(v)]$, which therefore
equals $F([0:v])$. For $a_0\ne0$, normalize by $a_0$; the
chart formula and multiplicativity give the homogeneous
formula. Its restriction to $[1:a:0:\cdots:0]$ determines
$\sigma(a)$, proving uniqueness.
\end{proof}
All chart identities are identities of normalized representatives.
Use \code{Projectivization.mk_eq_mk_iff} to convert equality of
points to a scalar multiple and then compare a coordinate equal
to one. This avoids any choice of compatible representatives.
The case $c=0$ above must be separated; one cannot invoke the
uniqueness of a line through two coincident points.

\node{C04}{HyperplaneGT.DiagramAut.toGalois}{I03--I04 and C03}
\begin{proposition}\label{prop:c04}
There is a group homomorphism
$\eta:\Aut_2\to(k\simeq_{k_0}k)$ whose underlying field
automorphism induces $\pi_g$.
\end{proposition}
\begin{proof}
The fixed coefficient points of $x_i=0,x_i=1,x_i-x_j=0$
and the deliberately fixed $o$ supply exactly the frame needed
in C03. For $s\in S$, the fixed hyperplane $x_i=s$ has
coefficient point $[-s:0:\cdots:1:\cdots:0]$. Apply the
homogeneous formula and compare the coordinate equal to one:
$\sigma(s)=s$. A field automorphism fixes the prime field.
Its fixed elements form a subfield, since they are closed under
the field operations. Therefore it fixes the subfield generated
by $S$ and the prime field, namely $k_0$. This packages it as
a $k_0$-algebra equivalence. Uniqueness in C03 and the
composition assertion in I01 prove the group law.
\end{proof}
For the generated-subfield step use
\code{RingHom.eqOn_field_closure} in
\code{Mathlib/Algebra/Field/Subfield/Basic.lean:505}, or use
subfield closure induction with the operations just listed.

\Needspace{14\baselineskip}
\section{The scalar kernel, splitting, and full graded algebra}

\node{S01}{HyperplaneGT.DiagramAut.scalar_of_toGalois_eq_one}{C04, I01--I02 and L01}
\begin{lemma}\label{lem:s01}
If $\eta(g)=1$, a unique unit $u\in R^\times$ acts as $g$ on
every degree-one object, including $T$.
\end{lemma}
\begin{proof}
I01 and C04 imply $g(e_H)=u_He_H$. The restriction of $g$ to
each generator line is an equivalence, so $u_H$ is a unit:
write its inverse on that line as multiplication by $v_H$ and
evaluate both composites at $e_H$ to get $u_Hv_H=1$.
If $H,K,L$ form a dependent triple, then $e_K-e_L\in U_H^1$.
Its image also lies there by I02's annihilator proof. Apply
the trace map; the result is
$(u_K-u_L)e_{H\cap K}=0$. Evaluation at that label gives
$u_K=u_L$. Permute the triple to equate all three multipliers.

Two distinct intersecting hyperplanes have independent affine
equations $f,g$ and independent normal parts by G02. The third
equation $f+g$ defines a different hyperplane through their
codimension-two intersection, so their multipliers agree. If
two distinct hyperplanes are parallel, choose a normal vector
independent of their common normal; the normal-vector space
has dimension $n\ge2$. The resulting hyperplane meets each,
and the preceding argument equates both multipliers to its
multiplier. Thus all $u_H$ equal one unit $u$. Finite basis
expansion gives $g=u\id$ on $M$. Its finite-stage restrictions
give the same equality by U02. There is a marking in some stage
by cofinality; its pullback from $T$ is injective, so naturality
gives $g_T=u\id$. Evaluation at any generator proves uniqueness.
\end{proof}

\node{S02}{HyperplaneGT.galoisGradedAut}{G04, G06 and the explicit OS presentation}
\begin{lemma}\label{lem:s02}
For each $\sigma\in G$, $e_H\mapsto e_{\sigma H}$ gives a
natural graded $R$-algebra automorphism $P_\sigma$ of the full
diagram. These automorphisms form an action, act identically
on $T$.
\end{lemma}
\begin{proof}
The permutation is defined at each stage because its arrangement
is $G$-stable. Its linear extension on the free generator module
extends to the exterior algebra by its universal property.
G04 preserves empty intersections and the codimensions of all
nonempty intersections; hence it carries every empty-intersection
word and every dependent boundary relator to another such
relator. It descends to the OS quotient. Use $\sigma^{-1}$ for
its inverse. Inclusion naturality and the group law hold on
each generator and hence on the algebra generated by them.
All marked equations have coefficients in $k_0$, so both of
their labels are fixed. Define the map at $T$ to be the identity;
agreement on its two generators proves marking naturality.
\end{proof}
This construction must check \emph{all} OS relators, not just the
quadratic ones. The quotient universal property already used in
\code{Presentation/Basic.lean} provides the implementation pattern.
Semilinear changes of the geometric coordinates give $R$-linear
maps of this presentation; they are not $k$-morphisms to which a
theory's ordinary pullback can be applied.

After C04, prove the separate comparison lemma
\code{HyperplaneGT.galoisGradedAut_toGalois}:
$\eta(P_\sigma)=\sigma$. Its proof is that the contragredient
of $P_\sigma$ sends $\delta_H$ to $\delta_{\sigma H}$;
C03's uniqueness therefore identifies its extracted field
automorphism with $\sigma$. Place this comparison in
\code{Reconstruction/LocalRing.lean}, not in the independent
\code{Diagram/GaloisAction.lean} module.

\node{S03}{HyperplaneGT.quadraticAutEquiv}{C04 and S01--S02}
\begin{theorem}\label{thm:s03}
For a nonzero commutative local ring there is a group equivalence
\[
 R^\times\times G\simeq\Aut_2,\qquad
 (u,\sigma)\longmapsto uP_\sigma.
\]
\end{theorem}
\begin{proof}
Multiplication by a unit preserves every quadratic relation span
and commutes with all diagram maps. It commutes with every
$R$-linear automorphism, so the displayed map is a homomorphism.
For $g$, let $\sigma=\eta(g)$. Then $gP_\sigma^{-1}$ is in the
kernel of $\eta$; S01 writes it uniquely as $u\id$. This proves
surjectivity. If $uP_\sigma=vP_\tau$, applying $\eta$ gives
$\sigma=\tau$, and evaluation at a generator gives $u=v$.
It also gives an explicit inverse and both inverse equations.
\end{proof}
Use a \code{MulEquiv}, not a bare bijection or an equivalence
of types. The scalar embedding is an injective group homomorphism.

\node{S04}{HyperplaneGT.gradedAutEquiv}{S02--S03 and full OS generation}
\begin{theorem}\label{thm:s04}
Restriction to degree one is a group equivalence
$\Aut_{\rm gr}\simeq\Aut_2$. Consequently both groups modulo
the central scalar subgroup are canonically $G$.
\end{theorem}
\begin{proof}
A graded algebra automorphism preserves the product kernel,
and its inverse proves equality of the kernels, so restriction
lands in $\Aut_2$. An $R$-algebra homomorphism of an OS algebra
is determined by its generators, because the exterior algebra
is generated by them and the quotient is onto. Equivalently
use \code{OSAlgebra.span_word_length} degree by degree and
multiplicativity. This proves injectivity of restriction.
By S03 every quadratic automorphism is $uP_\sigma$. On a word
of length $q$, multiplication by $u^q$ preserves the homogeneous
relator ideal; it gives a graded automorphism with inverse
defined using $u^{-1}$. Composing with S02 extends $uP_\sigma$,
proving surjectivity. These are group homomorphisms by their
generator formulas. For the quotient, the composite with
projection $R^\times\times G\to G$ is onto with kernel exactly
the scalar subgroup. The group quotient universal property and
the section $\sigma\mapsto[(1,\sigma)]$ give explicit mutually
inverse homomorphisms to $G$.
\end{proof}
No statement that the OS algebra is quadratic is used or true
in this generality. Extension becomes available only after the
entire marked diagram has forced the form $uP_\sigma$.

\node{S05}{HyperplaneGT.coefficientDomainAutEquiv}{A02--A03 and S02--S04 over a field}
\begin{theorem}\label{thm:s05}
The conclusions of S03--S04 hold over every integral domain $R$.
\end{theorem}
\begin{proof}
Let $K$ be its fraction field. Define the base-changed diagram
directly as the free $K$-module on each stage's labels, with the
same quadratic generators and the same marking matrices.
Extend the finite matrix of each component of $g$ along
$R\hookrightarrow K$. Extending the inverse matrix proves
invertibility, and extending the naturality matrix equations
proves naturality. For each quadratic relator $q$, its image
under $\bigwedge^2g$ is a finite $R$-linear combination of the
quadratic generators by A02 and the hypothesis on $g$. Extend
this equality along $R\to K$. Since those same generators
span the quadratic relation submodule over $K$, it proves
preservation over $K$; apply the inverse for equality. This
constructs $g_K\in\Aut_2(K)$ without an exactness assumption
on a new cohomology theory or on its scalar extension.

The field $K$ is local, so S03 gives $g_K=cP_\sigma$.
Choose a finite stable stage containing one hyperplane $H$;
it also contains $\sigma H$. The matrix entry at row
$\sigma H$, column $H$ is an element $r\in R$ with image
$c$. For the inverse matrix choose column $\sigma H$, row
$H$, obtaining $s\in R$ with image $c^{-1}$. Injectivity of
$R\to K$ gives $rs=1$ and $sr=1$, so $r$ defines a unit $u$.
Every entry of $g$ equals the corresponding entry of
$uP_\sigma$, because their images in $K$ agree. This applies
at all finite stages and at $T$. Thus $g=uP_\sigma$ over $R$.
Uniqueness follows over $K$, or by the same generator argument.
S02 constructs the graded extension, and S04's generation proof
is valid over every coefficient ring.
\end{proof}
\emph{Implementation.} Use the explicit Finsupp bases to define
matrix scalar extension. This avoids needing any new theorem on
localization of cohomology kernels. Existing
\code{IsFractionRing.injective}, in
\code{Mathlib/RingTheory/Localization/FractionRing.lean:144},
justifies every descent of an equality. If a full algebra base-change
comparison is desired, use \code{OSAlgebra.baseChangeEquiv} and
\code{OSAlgebra.map_baseChangeEquiv} in
\code{Presentation/BaseChange.lean:116,146}; their generator
formulas identify the same construction.

\Needspace{14\baselineskip}
\section{Transport to the theory and examples}

\node{E01}{HyperplaneGT.OrdinaryTheory.reconstruction}{A04 and S03--S05}
\begin{theorem}\label{thm:e01}
For an existing geometric ordinary theory satisfying C1--C5,
the marked diagonal cohomology diagram has the group equivalences
stated in Section~\ref{sec:boundary}. They carry $(u,\sigma)$
to $[H]\mapsto u[\sigma H]$ and act by $u$ on degree one of $T$.
\end{theorem}
\begin{proof}
A04 identifies the entire marked cohomology diagram, its
degree-one modules, and its full degree-two kernels with the
explicit OS diagram. Conjugate every component of an
automorphism by these equivalences. Naturality and preservation
of kernels are preserved under conjugation, and the inverse
conjugation proves a group equivalence of automorphism groups.
The same argument with the diagonal algebra equivalences gives
the graded group equivalence. Their generator formulas identify
scalars and the action in the statement. Apply S03--S05.
\end{proof}
This is the only point, after A04, where cohomology enters the
reconstruction. Keep the preceding statements independent of
theory universes, grading implementations and divisor pairs.

\node{E02}{HyperplaneGT.cofinalRestrictionEquiv}{E01 and G06}
\begin{corollary}\label{cor:e02}
Restriction to any cofinal subsystem of the chosen stable
arrangements induces an isomorphism of either automorphism
group. Two eligible realizations with the same coefficient
ring have canonically isomorphic marked diagonal diagrams.
\end{corollary}
\begin{proof}
For the first assertion the restriction homomorphism carries
the family $uP_\sigma$ to the family with the same formula.
Both groups are classified by E01, so the map is the identity
on $R^\times\times G$ and is an equivalence. Do not claim that
arbitrary diagram automorphisms extend from a cofinal subsystem
before proving reconstruction.

For the second assertion, at every stage compose one theory's
canonical OS equivalence with the inverse of the other's, in
the appropriate order. A04 makes these maps natural for every
marking and inclusion. They send $[H]$ to $[H]$, and generation
proves uniqueness. They therefore identify the classified
actions, independently of any comparison theorem between the
two realizations.
\end{proof}

\node{E03}{HyperplaneGT.deRham_reconstruction}{E01 and the constructed de Rham instance}
\begin{corollary}\label{cor:e03}
If $\operatorname{char}k=0$, the ordinary algebraic de Rham
cohomology diagram of the complements has
\[
 \Aut_{\rm gr}\simeq\Aut_2\simeq
 k^\times\times\Gal(k/k_0),
 \qquad [df_H/f_H]\longmapsto u[df_{\sigma H}/f_{\sigma H}].
\]
\end{corollary}
\begin{proof}
Take \code{Geometric.DeRham.ordinaryTheory}. Its coefficients
are the Laurent weight extension of $k$, with weight-zero
coefficient ring $k$. The existing diagonal and degree
identifications in \code{DeRham/Schemes/Corollaries.lean}
identify its diagonal with the usual graded de Rham algebra
and its logarithmic class with $df/f$. A field is a nonzero
local ring. Apply E01 and transport through these equivalences.
\end{proof}
This action is $k$-linear. The usual transport of differential
forms by $\sigma$ is $\sigma$-semilinear on scalar coefficients.
The theorem constructs the former from the logarithmic
presentation and must not label it the latter.

For singular cohomology, the desired future input is the
Laurent weight package
$E^{p,w}(X)=H^p(X_\iota(\C),R)t^w$ for a fixed embedding
$\iota:k\hookrightarrow\C$, with complex-oriented residues.
For geometric \'etale cohomology the desired input is
$E^{p,w}(X)=H^p_{\mathrm{\acute et}}(X,R(w))$,
$R=\Z/\ell^m,\Z_\ell$ or $\Q_\ell$ and
$\ell\ne\operatorname{char}k$. Each is an instance of E01
\emph{once} its C1--C5 data have been constructed. Neither is
an imported theorem or a new Lean axiom in this blueprint.
Their ordinary mathematical realizations and verification of
the hypotheses are discussed in the paper. Implementing their
foundations is separate from formalizing this new reconstruction
argument. In particular the blueprint's unconditional named
realization theorem is E03, using only the existing instance.

\Needspace{14\baselineskip}
\section{Lean organization, dependency order, and completion criteria}

The following module names are proposals under
\code{OrlikSolomon/HyperplaneGT/}; every public declaration should
document its node ID. They do not replace the existing
\code{docs/ARCHITECTURE.md}. Split files further if they approach
the repository's 800-line limit.
\small
\begin{longtable}{@{}p{.30\textwidth}p{.13\textwidth}p{.50\textwidth}@{}}
\toprule Proposed module & Nodes & Output and immediate dependencies\\\midrule\endhead
\code{Geometry/Coordinates.lean}&G01--G03&Affine equation, graph, trace and incidence API; existing hyperplane geometry.\\
\code{Geometry/Galois.lean}&G04--G06&Coefficient action, finite orbits, stable finite systems and markings.\\
\code{Algebra/LowDegrees.lean}&A01--A02&Canonical generator basis and exact quadratic kernel; existing presentation and grading.\\
\code{Diagram/Basic.lean}&A03&Explicit finite quadratic and graded diagrams; automorphism groups.\\
\code{Diagram/Universal.lean}&U01--U02&Free universal module, finite-stage embeddings and induced equivalences.\\
\code{Algebra/Dual.lean}&U03--U05&Product coordinates, commuting, primitive vectors over local rings.\\
\code{Reconstruction/Trace.lean}&L01--L03&Trace kernel, inertia, decomposition and coordinate restrictions.\\
\code{Reconstruction/Local.lean}&L04--L05&Local transfer calculation and unit-supported detection.\\
\code{Reconstruction/Projections.lean}&P01--P02&Intrinsic coordinate-span recognition and invariance.\\
\code{Reconstruction/Incidence.lean}&I01--I04&Hyperplane permutation, collineation and fixed markings.\\
\code{Geometry/Collineation.lean}&C01--C03&Pure coordinate field-reconstruction theorem, independent of cohomology and $R$.\\
\code{Reconstruction/Galois.lean}&C04, S01&Extracted Galois homomorphism and scalar kernel.\\
\code{Diagram/GaloisAction.lean}&S02&Action on full OS algebras; independent of local detection.\\
\code{Reconstruction/LocalRing.lean}&S03--S04&Local coefficient classification and extension to the graded algebra.\\
\code{Reconstruction/Domain.lean}&S05&Matrix extension to a fraction field and descent.\\
\code{Cohomology/Adapter.lean}&A04&Logarithm naturality and finite diagram comparison; existing ordinary theory.\\
\code{Cohomology/Main.lean}&E01--E02&Theory-level classification and cofinal/realization independence.\\
\code{Cohomology/DeRham.lean}&E03&Application to the existing de Rham construction.\\\bottomrule
\end{longtable}
\normalsize

The main dependency chain is
\[
 \begin{gathered}
 (G01\text{--}G06,A01\text{--}A03)\longrightarrow
 U01\text{--}U05\longrightarrow L01\text{--}L05\longrightarrow
 P01\text{--}P02\longrightarrow I01\text{--}I04,\\
 (C01\longrightarrow C02\longrightarrow C03,\ I03\text{--}I04)
 \longrightarrow C04\longrightarrow S01,
 \qquad (S01,S02)\longrightarrow S03\text{--}S05,\\
 (A04,S03\text{--}S05)\longrightarrow E01\longrightarrow E02\text{--}E03.
 \end{gathered}
\]
The pure C01--C03 theorem should take an arbitrary field and a
line-preserving permutation with the stated fixed frame; its
proof does not depend on the construction of I03. Similarly
U03's dual-coordinate equivalence and U05's local-ring lemma
should be stated for an arbitrary index type, not hyperplanes.
S02 is built directly from G04 and the OS presentation, so it
does not depend on the reconstructed homomorphism; only its
last identification with $\eta$ depends on C04.

\subsection{Suggested public interfaces}

The following is schematic Lean, not checked code. Names of
existing declarations retain their current namespaces; proposed
structures must not add geometric or cohomological conclusions
as fields. In particular \code{QuadraticDiagram} stores the
explicit free modules and relation spans of A03; it is not an
axiom package asserting reconstruction.
\par\smallskip
\noindent\begin{minipage}{\linewidth}
\begin{lstlisting}
-- Pure geometry: arbitrary field k, finite dimension n >= 2.
-- Its conclusion is a uniquely determined ring equivalence of k.
HyperplaneGT.coordinateFieldEquiv

-- The core diagram is explicit, so no cohomology is assumed.
-- R is a nonzero commutative local ring; F is cofinal and G-stable.
HyperplaneGT.quadraticAutEquiv :
  MulEquiv (Prod (Units R) (AlgEquiv k0 k k))
    (QuadraticDiagram.Aut R F S)

-- Same conclusion with R an integral domain.
HyperplaneGT.coefficientDomainAutEquiv

-- T is an existing Geometric.OrdinaryTheory, not a new weaker theory.
HyperplaneGT.OrdinaryTheory.reconstruction

-- Unconditional in characteristic zero via the existing instance.
HyperplaneGT.deRham_reconstruction
\end{lstlisting}
\end{minipage}
\par\smallskip
The actual declarations will quantify the fields, algebra maps,
finite dimension, coefficient ring, generating set and cofinal
system explicitly. Use \texttt{[Nontrivial R]} with
\texttt{[IsLocalRing R]}, and \texttt{[IsDomain R]} in the second
version. The group $(k\simeq_{k_0}k)$ exists for any field
extension; require algebraicity for G05, and the perfect-field
algebraic-closure hypotheses at the public geometric theorem
matching the paper. No theorem in the pure algebra layer should
carry these latter unused assumptions. Give the public
equivalences \code{apply_generator}, \code{apply_marked},
\code{symm_apply_toGalois} and scalar-kernel lemmas.

\subsection{Regression statements and mathematical limits}

These are mathematically meaningful checks of the generality,
not additional hypotheses to help a proof succeed.
\begin{enumerate}
\item \textbf{Characteristic two and nilpotents.} Instantiate the
local coefficient theorem at $\F_2$ and at
$\F_2[\epsilon]/(\epsilon^2)$. Exterior squares are alternating
by the existing OS theorem; no step divides by $2$. The only
division in L04 is by a coordinate proved to be a unit. In
I04, distinguish the two-label permutation combinatorially.
The residue-field criterion alone must never replace the
surjectivity definition of primitive over an infinite product.
\item \textbf{Dimension.} The new theorem assumes $n\ge2$,
exactly where L05 needs an intersecting graph, I01 a second
coordinate marking, C03 a coordinate plane, and S01 a transverse
normal. Do not weaken this to $n\ge1$. Degree computations A01--A02
remain valid in dimensions zero and one and for the empty
arrangement, with their evident empty bases.
\item \textbf{Zero coefficient ring.} All purely algebraic OS,
trace-map and diagram definitions can be made over the zero ring.
The reconstruction theorem cannot: then every module and every
automorphism group is trivial. The hypotheses
\texttt{[Nontrivial R]} enter U05, uniqueness of inertia labels,
and detection of distinct trace coordinate lines. Do not add
nontriviality to the earlier structure theorem.
\item \textbf{Disconnected coefficients.} Over $R=F\times F$,
the idempotents split every finite diagram into two copies over
$F$. Choose distinct $\sigma,\tau\in G$. The action
$(P_\sigma,P_\tau)$ preserves every full OS relation and every
marking. It cannot equal a common $uP_\rho$, since projecting to
each factor and comparing generator labels forces
$\rho=\sigma=\tau$. Thus a theorem over arbitrary commutative
rings is false whenever $G$ has two elements. For
$k_0=\Q,k=\overline\Q$ such automorphisms exist, for example
by extending the two roots of $X^2-2$ to the algebraic closure
using Mathlib's normal-extension lifting API. No such example
needs to be imported for the theorem; the conditional
counterexample just proved already identifies the obstruction.
\item \textbf{Geometric points.} The argument uses affine
hyperplanes over $k$ and actual smooth scheme complements for
cohomology. It never requires a rational point of a complement
or a descended complement over $k_0$. Existing structure-theorem
coverage of complements without rational points must be retained.
\item \textbf{Finite stage compatibility.} A relation witness,
a marking, or a selected list of hyperplanes may be moved to a
common stage using G06. Do not postulate that the cofinal
family is closed under unions. U02 must prove independence of
stage and preserve the full quadratic relation submodule.
\item \textbf{No extra naturality.} Recover coordinate subspace
spans by P01--P02; do not add arbitrary projections as diagram
arrows. There must be one common target $T$ for all markings;
separate targets would invalidate I04's common permutation.
\item \textbf{No dual or quadratic shortcuts.} The dual is the
full algebraic product, with maps induced from the free module.
Use the explicit annihilator and trace calculations rather than
an unproved reflexivity assertion. The graded extension must
come after S03. S05 can use matrices and relation spans rather
than unproved localization theorems about an arbitrary theory.
\end{enumerate}

Completion means that all proposed modules elaborate without
new axioms or placeholders, the generator formulas and group
equivalence laws are proved, and the theorem is instantiated
with the existing de Rham theory. Existing repository build and
review rules apply at that implementation stage. This document
itself is a mathematical specification; it does not
claim that these new declarations have already been formalized.
\end{document}
